\documentclass[10pt]{article}
\usepackage[T1]{fontenc}
\usepackage[utf8]{inputenc}
\usepackage[french]{babel}
\usepackage{amsmath,amssymb,mathrsfs,mathpazo}
\usepackage{geometry}
\geometry{a4paper,margin=2cm}
\usepackage{graphicx,array,tabularx,tabularray,colortbl,multirow,multicol,booktabs,enumitem,pifont,ulem,fancybox,diagbox,cancel}
\usepackage{pstricks,pstricks-add,pst-plot,pst-tree,pst-node,pst-text,pst-func,pst-math,pst-bspline,pst-3dplot,pst-eucl,pst-solides3d}
\let\pscalculate\undefined
\usepackage{pst-bezier}
\usepackage{tikz,pgfplots,tkz-tab}
\usetikzlibrary{arrows,arrows.meta,patterns,calc,positioning,trees,decorations.markings,shapes.geometric,shapes.misc,shapes.arrows,intersections,decorations.pathmorphing,decorations.pathreplacing,backgrounds,shadows,babel}
\pgfplotsset{compat=1.18}
\usepgfplotslibrary{fillbetween}
\usepackage[np]{numprint}
\usepackage{listings,makecell,pas-tableur,twoopt}
\let\barre\undefined
\usepackage[pstricks]{tablvar}
\definecolor{darkgreen}{RGB}{0,100,0}
\providecommand{\red}{\color{red}}
\providecommand{\blue}{\color{blue}}
\providecommand{\green}{\color{green}}
\providecommand{\black}{\color{black}}
\providecommand{\white}{\color{white}}
\providecommand{\gray}{\color{gray}}
\providecommand{\cyan}{\color{cyan}}
\providecommand{\magenta}{\color{magenta}}
\providecommand{\yellow}{\color{yellow}}
\providecommand{\violet}{\color{violet}}
\providecommand{\orange}{\color{orange}}
\providecommand{\euro}{\texteuro}
\providecommand{\EUR}{\texteuro}
\providecommand{\ieme}{\textsuperscript{e}}
\providecommand{\degres}{\ensuremath{^\circ}}
\providecommand{\degree}{\ensuremath{^\circ}}
\providecommand{\vv}[1]{\overrightarrow{#1}}
\providecommand{\Coord}[1]{#1}
\providecommand{\decosix}{}
\providecommand{\starredbullet}{\ensuremath{\bullet}}
\pagestyle{empty}
\setlength{\parindent}{0pt}
\usepackage{adjustbox}\DeclareUnicodeCharacter{2236}{:}\DeclareUnicodeCharacter{2190}{\ensuremath{\leftarrow}}
\providecommand{\R}{\mathbb{R}}
\providecommand{\N}{\mathbb{N}}
\providecommand{\Z}{\mathbb{Z}}
\providecommand{\Q}{\mathbb{Q}}
\providecommand{\D}{\mathbb{D}}
\providecommand{\C}{\mathbb{C}}
\providecommand{\e}{{\mathrm{e}}}
\providecommand{\vect}[1]{\overrightarrow{#1}}
\providecommand{\vectt}[1]{\overrightarrow{\mathrm{#1}}}
\providecommand{\vvt}[1]{\overrightarrow{\mathrm{#1}}}
\providecommand{\np}[2][]{{#2~#1}}
\providecommand{\vv}[1]{{\overrightarrow{#1}}}
\renewcommand{\d}{{\mathrm{d}}}
\providecommand{\degres}{{^\circ}}
\providecommand{\degre}{{^\circ}}
\providecommand{\blue}{\color{blue}}
\providecommand{\red}{\color{red}}
\providecommand{\enskip}{\quad}
\providecommand{\llbracket}{\lbrack\!\lbrack}
\providecommand{\rrbracket}{\rbrack\!\rbrack}
\providecommand{\euro}{{\text{€}}}
\providecommand{\no}{{\text{n°}}}
\providecommand{\interval}[2]{{[#1;#2]}}
\providecommand{\intervFO}[2]{{[#1;#2[}}
\providecommand{\dotfill}{\ldots}

\providecommand{\ds}{\displaystyle}
\providecommand{\textding}[1]{#1}
\providecommand{\Oij}{\(\mathrm{O};\vec{\imath},\vec{\jmath}\)}
\providecommand{\Oijk}{\(\mathrm{O};\vec{\imath},\vec{\jmath},\vec{k}\)}
\providecommand{\Ouv}{\(\mathrm{O};\vec{u},\vec{v}\)}
\providecommand{\intervFF}[2]{\ensuremath{[#1;#2]}}
\renewcommand{\euro}{\ifmmode\text{\texteuro}\else\texteuro\fi}
\DeclareUnicodeCharacter{00B0}{\ifmmode^{\circ}\else\textdegree\fi}
\DeclareUnicodeCharacter{20AC}{\ifmmode\text{\texteuro}\else\texteuro\fi}
% Énoncés : Éducation nationale. Transcription et corrigés : APMEP, Denis Vergès et contributeurs. Découpage et classement : CoopMaths, Sébastien Lozano, Éric Elter et contributeurs. Adaptation de la présentation : Maths Explorer.
% CC BY-SA : https://coopmaths.fr/www/about/
% https://coopmaths.fr/alea/static/eam/2026/tex/eam_technologique_2026_06_etrangers_automatismes_cor.tex
\begin{document}
\noindent\textit{APMEP / CoopMaths. CC BY-SA. Présentation adaptée par Maths Explorer.}\par\medskip
\subsubsection*{Question 1}
On écrit $5=\dfrac{10}{2}$. Ainsi :
\[5-\dfrac{3}{2}=\dfrac{10}{2}-\dfrac{3}{2}=\dfrac{7}{2}.\]

La bonne réponse est \textbf{b.}

\subsubsection*{Question 2}
On utilise la propriété $(a^m)^n=a^{m \times n}$ :
\[\left(2^3\right)^{-5}=2^{3\times(-5)}=2^{-15}.\]

La bonne réponse est \textbf{b.}

\subsubsection*{Question 3}
Une baisse de $20\,\%$ correspond à un coefficient multiplicateur égal à : $1-\dfrac{20}{100}=0{,}8$.

Le nouveau prix est donc $990\times 0{,}8$.

La bonne réponse est \textbf{d.}

\subsubsection*{Question 4}
Obtenir un nombre supérieur ou égal à 5 signifie obtenir 5 ou 6. Donc, dans le cas présent :
\[P(5 \cup 6)=P(5)+P(6)=0{,}3+0{,}3=0{,}6.\]

La bonne réponse est \textbf{c.}

\subsubsection*{Question 5}
Un produit est nul si, et seulement si, l'un de ses facteurs est nul :
\[x-2=0 \quad \text{ou} \quad 2x+1=0 \Rightarrow 2x=-1 \Rightarrow x=-\dfrac{1}{2}.\]
On obtient donc $x=2$ ou $x=-\dfrac{1}{2}$.

La bonne réponse est \textbf{a.}

\subsubsection*{Question 6}
Il y a $5+7+8+10=30$ élèves. La médiane est donc la moyenne des 15\ieme{} et 16\ieme{} valeurs de la série ordonnée.

Les notes 7 occupent les rangs 1 à 5, les notes 10 les rangs 6 à 12, les notes 12 les rangs 13 à 20. Les 15\ieme{} et 16\ieme{} valeurs valent donc toutes les deux 12. La médiane est donc $12$.

La bonne réponse est \textbf{a.}

\subsubsection*{Question 7}
Sur la courbe, la fonction est positive ou nulle lorsque la courbe est située sur ou au-dessus de l'axe des abscisses.

On lit que la courbe coupe l'axe des abscisses en $x=-3$ et en $x=-1$. Elle est au-dessus de l'axe sur $[-4\,;\,-3]$ puis sur $[-1\,;\,3]$. L'ensemble des solutions est donc : $[-4\,;\,-3]\cup[-1\,;\,3]$.

La bonne réponse est \textbf{c.}

\subsubsection*{Question 8}
La proportion d'élèves pratiquant un sport est : $\dfrac{150}{500}=0{,}3=30\,\%$.

La bonne réponse est \textbf{c.}




\begin{center}
\end{center}

\end{document}
